\cotefile{71}
\scanpage{36}\mspage{$\underline{\overline{\mathrm{IX}}}$}
Comme $\Gamma$ est un $J$-polygone et $\partial E(\Gamma)$ est un polygone convexe il est possible de trouver un disque $D_s$ qui vérifie b) de la def.~1 par rapport à $\Gamma$ et à $\partial E(\Gamma)$. Par le lemme $\underline{\overline{\mathrm{IX}}}$ on sait que l'angle $\alpha_1$ du secteur angulaire interieur (de $D_s$) determiné par $\partial E(\Gamma)$ est inferieur à $\pi$, \deleted{alors nécessairement $\alpha_1<\pi$}, car $s$, $s'$ et $s''$ ne sont pas colineaires. Par conséquent l'angle $\alpha_2$ qui sous-tend le secteur angulaire exterieur vérifie $\alpha_2>\pi$.

Mais $\operatorname{ext}\partial E(\Gamma)\subset\operatorname{ext}\Gamma$ implique qu'on peut trouver dans $D_s$ un angle $\alpha'_2>\pi$ qui sous-tend un secteur angulaire contenu dans $\operatorname{ext}\Gamma$, par conséquent le secteur angulaire interieur de $D_s$ \deleted{par} donné par $\Gamma$ a un angle $<\pi$.

Mais tout polygone convexe a au moins trois sommets et cela suffit à prouver le lemme.

On est en mésure de montrer le :

\uline{Lemme X}I : Si $\Gamma$ est un $J$-polygone \inserted{à $n\unsure{>|\geq}3$ \unsure{côtés|cotés}} il existe un couple de sommets $s,s'\in S(\Gamma)$ tels que $s'$ est visible de $s$ de l'interieur de $\Gamma$.
% Pencil, slanted in the left margin beside the statement.
\margin{faux pour $n=3$!}

% Pencil line below the statement.
\margin{OPS $\Gamma$ non convexe (sinon c'est clair grâce à --)}

\uline{Démonstration} : Soit $s$ un sommet satisfaisant la propriété ennoncée au lemme $\underline{\overline{\mathrm{VIII}}}$%
% Dark-ink correction, left margin beside the reference.
\margin{\uline{lemme X}}, soient $s'$, $s''$ les deux sommets \deleted{consecutifs} \inserted{\unsure{adjacents}} à $s$ et $T$ le tri\uline{an}gle%
% Dark-ink note beside the word triangle and its short underline.
\margin{ambigu} de sommets $s,s',s''$. Si $(T-([s,s']\cup[s,s'']))\cap\Gamma=\unsure{\phi|\varnothing}$ on a que $]s',s''[\cap\Gamma=\phi$, donc $]s',s''[\subset\operatorname{int}\Gamma$ ou bien $]s',s''[\subset\operatorname{ext}\Gamma$. Mais $\inserted{]s',s''[\subset}T-([s,s']\cup[s,s''])\subset\deleted{\unsure{\operatorname{ext}}}\operatorname{int}\Gamma$ (puisqu'il est connexe et et son intersection avec $\operatorname{int}\Gamma$ est non vide par le choix de $s$) donc $]s',s''[\subset\operatorname{int}\Gamma$.
% Dark-ink faux, with a lighter curved leader below it to the following pencil note.
\margin{\tikz[baseline=(fx.base)]{\node[inner sep=1pt] (fx) {\uline{faux}};\draw[gray] ([xshift=.15cm]fx.south) to[out=-100,in=-170] ++(1.6,-.12);}}
% Pencil note extends across the space below the first conclusion.
\margin{il faut prendre $s$ comme dans \unsure{la|le} \uline{dém} du lemme $\underline{\overline{\mathrm{VIII}}}$ (pas son énoncé!)}

% Dark-ink marginal words cancelled by repeated oblique hatching; exclamation survives.
\margin{\deleted{\illegible{2}}!}
% The intact inclusion is preceded by a separate long diagonal leader with a lower hooked arrowhead; it crosses the underline to its left.
\uline{Si} \uline{$(T-([s,s']\cup[s,s'']))\cap\Gamma\mathrel{\text{\tikz[baseline=(rel.base)]{\node[inner sep=0pt] (rel) {$\subset$};\draw[overlay,->] (-.18,.36) -- (-.48,-.42);}}}]s',s''[$} on doit pouvoir trouver au moins un sommet $z\in]s',s''[$ (puisque $\Gamma\cap]s',s''[\ne]s',s''[$). On a aussi $\mathring{T}\cap\Gamma=\phi$ donc $]s,z[\subset\text{\tikz[baseline=(ic.base)]{\node[inner xsep=3pt,inner ysep=2pt] (ic) {$\mathring{T}\subset\operatorname{int}\Gamma$};\draw (ic.center) ellipse[x radius=1.02cm,y radius=.40cm];\draw ([yshift=-.40cm]ic.center) to[bend right=23] ++(.25,-.24);}}$ (puisque $\mathring{T}\subset\mathbb{R}^2-\Gamma$ et $\mathring{T}\cap\operatorname{int}\Gamma\ne\phi$)
% Circled lower-left note. The first mark may be compressed est/mais or a short arrow.
\margin{\tikz[baseline=(m.base)]{\node[inner sep=1pt] (m) {\shortstack{\unsure{est|\ensuremath{\rightarrow}|mais} non\\vide?}};\draw (m.center) ellipse[x radius=1.24cm,y radius=.47cm];\draw (-1.14,-.12) -- (-.90,.60);}}
% Pencil note at the lower right, linked to the circled inclusion above.
\margin{dire avant! (et tu \unsure{as au|auras|as un} lemme $\underline{\overline{\mathrm{VIII}}}$)}
% archive sheet number: 35
